\chapter[Aprendizados e Perspectivas]{Aprendizados e Perspectivas}
\label{cap:conclusao}

Este capítulo apresenta uma reflexão crítica sobre o desenvolvimento do Observatório da Jornada do Paciente Oncológico, articulando os resultados técnicos aos objetivos do Relatório Técnico Profissional, às limitações enfrentadas e às perspectivas de continuidade.

\section{Aprendizado Individual e Coletivo}

O desenvolvimento da solução exigiu a integração de competências típicas de Sistemas de Informação: engenharia de software front-end, tratamento de dados em volume e design de interfaces para comunicação pública. Em especial, destacam-se:

\begin{itemize}
	\item a compreensão de que dados oficiais, por si só, não garantem transparência --- é necessário oferecer meios de exploração compreensíveis;
	\item a prática de separar preparação de dados (pipeline offline em Node.js) e apresentação (React + Vite), decisão decisiva para viabilizar o uso no navegador diante de mais de um milhão de registros;
	\item o domínio de bibliotecas de visualização (Recharts e d3-geo) aplicadas a um problema real de saúde pública;
	\item a atenção à linguagem da interface: disponibilizar recortes sociodemográficos e geográficos de modo a capacitar o usuário a interpretar os indicadores.
\end{itemize}

Do ponto de vista acadêmico, o trabalho reforçou a articulação entre programação, visualização de informação e responsabilidade ética na comunicação de indicadores de saúde. Coletivamente, a experiência evidenciou o valor de documentar limitações metodológicas junto ao próprio produto, e não apenas no relatório.

\section{Desafios}

Entre os principais desafios enfrentados, e as formas de enfrentá-los, destacam-se:

\begin{itemize}
	\item \textbf{Volume e codificação do \acs{CSV}:} o arquivo do \acs{RHC} (latin1, centenas de megabytes) inviabilizou o parsing no cliente. A solução foi o script \texttt{preprocess.js} com leitura em streaming e geração de \acs{JSON} agregados.
	\item \textbf{Qualidade variável dos campos:} datas inválidas e categorias ``ignorado'' reduziram a cobertura de alguns indicadores. Mitigou-se o problema com regras explícitas de exclusão e com a exposição, na interface, de percentuais de informação ignorada.
	\item \textbf{Equilíbrio editorial:} foi necessário revisar textos da \acs{UI} e do relatório para privilegiar a exploração dos dados, sem induzir o leitor a conclusões pré-estabelecidas.
	\item \textbf{Ausência de suíte automatizada:} compensada por validação funcional manual, inspeção dos agregados e verificação do build Vite a cada entrega relevante.
\end{itemize}

Esses obstáculos reforçaram a importância de tratar a qualidade dos dados e a clareza metodológica como parte do produto entregue ao público.

\section{Avaliação Crítica}

Em relação aos objetivos propostos, a solução entrega uma plataforma pública capaz de:

\begin{itemize}
	\item apresentar indicadores globais de tempo entre diagnóstico e início do tratamento;
	\item permitir comparação descritiva com o prazo legal de sessenta dias \cite{lei12732};
	\item oferecer exploração geográfica e por variáveis, bem como filtros de perfil;
	\item operar sem backend, facilitando publicação para governo e sociedade.
\end{itemize}

Pontos fortes incluem a arquitetura estática, o pipeline reproduzível e a presença de seções de metodologia na própria aplicação. Pontos a melhorar, em uma eventual retomada do ciclo, seriam a introdução precoce de testes automatizados sobre as funções de agregação e a geração sistemática de capturas de tela para a documentação. Ainda assim, a arquitetura escolhida mostrou-se adequada ao escopo e ao propósito de transparência.

\section{Articulação com o Curso ou Formação}

O trabalho dialoga diretamente com disciplinas e competências do curso de Sistemas de Informação, tais como desenvolvimento web, modelagem e tratamento de dados (aqui adaptados a agregados), engenharia de software e interação humano-computador. A experiência evidenciou, ainda, a relevância de formação complementar em visualização de dados e em leitura crítica de indicadores de saúde pública --- temas que atravessam a fronteira entre computação e políticas públicas.

Uma lacuna percebida foi a necessidade de maior formalização de testes automatizados e de documentação de \acs{API} de dados, conteúdos que poderiam ser aprofundados em disciplinas de qualidade de software e de integração de sistemas.

\section{Limitações}

Além dos desafios já citados, reconhecem-se limitações estruturais da solução e da base:

\begin{itemize}
	\item a base \acs{IRHC} é hospitalar e não cobre, necessariamente, todos os casos oncológicos do país \cite{inca_irhc};
	\item diferenças entre grupos podem refletir composição de tumores, estadiamento e outros fatores de confusão, e não relações causais --- cabendo ao usuário a interpretação;
	\item a aplicação atual não oferece drill-down municipal nem \acs{API} programática;
	\item dependência de reprocessamento offline sempre que a base \acs{CSV} for atualizada;
	\item a validação concentrou-se em testes manuais, sem cobertura automatizada contínua.
\end{itemize}

\section{Perspectivas Futuras}

Como desdobramentos possíveis da plataforma, sugerem-se:

\begin{itemize}
	\item aprofundamento geográfico (município ou macrorregião), mantendo regras claras de amostra mínima;
	\item exposição de uma \acs{API} ou de arquivos versionados para reuso por gestores e pesquisadores;
	\item empacotamento como \acs{PWA} para uso offline de agregados já baixados;
	\item suíte de testes automatizados no pipeline de dados (unitários nas funções de agregação e regressão dos totais);
	\item painel de metadados da base (versão do \acs{CSV}, data de processamento, dicionário de campos);
	\item estudos colaborativos com especialistas em epidemiologia para enriquecer as notas metodológicas, preservando o caráter exploratório da ferramenta.
\end{itemize}

Em síntese, o Observatório cumpre o papel de tornar dados públicos mais acessíveis à análise por governo e cidadãos, contribuindo para a transparência e para a autonomia interpretativa do usuário.
